\documentclass[11pt]{article}

\usepackage[margin=1in]{geometry}
\usepackage{amsmath,amssymb,amsthm}
\usepackage{mathpartir}
\usepackage{underscore}   % lets identifier_names appear in \texttt{} without escaping
\usepackage{hyperref}
\usepackage{enumitem}
\usepackage{booktabs}

\newcommand{\inN}{\!\in\!}
\newcommand{\file}[1]{\texttt{#1}}
\newcommand{\op}[1]{\texttt{#1}}

\theoremstyle{definition}
\newtheorem{definition}{Definition}[section]
\newtheorem{lemma}[definition]{Lemma}
\newtheorem{theorem}[definition]{Theorem}
\newtheorem{corollary}[definition]{Corollary}
\newtheorem{remark}[definition]{Remark}
\newtheorem{pitfall}[definition]{Pitfall}

\title{Proof notes: what each result actually says and why it's true}
\author{Hoang Nguyen}
\date{}

\begin{document}
\maketitle

\begin{abstract}
Most theorems below sound obviously true: merging regions can't duplicate an id, a field write
somewhere active can't disturb a suspended region, a closed region only opens through its front
door. Intuition gives the headline, not the proof: each either hides an assumption that needs
deriving from the model's invariants, or is wrong until qualified in a way you'd only notice by
trying to break it. This is for someone who trusts the model (\file{writeup/main.tex},
\file{lean/gc/README.md}) but hasn't seen how these results actually got proved.

Each result: the claim, what intuition covers, where the real work is, how the proof gets there. A
lemma catalog and known pitfalls follow at the end, for lookup once you know which proof you want.

Scope: \file{Gc/Model/} and \file{Gc/Reachability/Reachable/}. \file{Gc/Reachability/Referencable/}
is excluded (deprecated, superseded by \file{Reachable/}; see \file{CLAUDE.md}).
\end{abstract}

\tableofcontents

\section{Preliminaries}

\subsection{The model}

Configuration $\langle S, H\rangle$ (stack, heap) and the nine operations' rules: full definitions
in \file{writeup/main.tex}'s ``Model''/``Operational Semantics''.\footnote{Compiles standalone, so
section names are given in prose, not as cross-file \texttt{\textbackslash ref}s.} Notation:
$\mathsf{RId}(rid)$/$\mathsf{OId}(oid)$ for references, $a.loc(\langle S,H\rangle)$ for where a
reference resolves ($\mathsf{Stk}(fid)$, $\mathsf{Rgn}(rid)$, undefined), $a.objAt(\langle
S,H\rangle)$ for the object an $\mathsf{OId}$ names.

\subsection{The ten invariants of \textit{ValidConfig}}
\label{sec:invariants}

\file{main.tex} defines all ten and bundles them into \textit{ValidConfig}. Not restated here, just
what each forbids:

\begin{center}
\begin{tabular}{@{}ll@{}}
\toprule
Invariant & What it rules out \\
\midrule
\textit{L1} & the same object id allocated twice \\
\textit{L2} & a frame executing inside a \textit{Closed} region \\
\textit{H1} & a region whose bridge object isn't in its own \textit{objMap} \\
\textit{H2} & two live $\mathsf{RId}(rid)$ references to the same region \\
\textit{H3} & a region's contents pointing outside that region \\
\textit{S1} & two stack frames sharing a region \\
\textit{S2} & a frame referencing a stack object owned by a later frame \\
\textit{S3} & a frame referencing a region owned only by a later frame \\
\textit{HS1} & a dangling $\mathsf{OId}$ \\
\textit{HS2} & a dangling $\mathsf{RId}$ \\
\bottomrule
\end{tabular}
\end{center}

\begin{remark}
\textit{HS1}/\textit{HS2} don't follow from the rest: \textit{H2}/\textit{H3}/\textit{S2}/\textit{S3}
all presuppose a reference that already resolves somewhere, none says anything about one that
resolves to nothing. Without them a stale reference could sit inert and later collide, uncaught,
with a freshly allocated id. (Surfaced while proving \textit{S2}/\textit{S3} for the three
\op{make*} operations, whose \op{freshObjectId}/\op{freshRegionId} only look at ids used as keys.)
\end{remark}

Four corollaries (\file{Validity.lean}), cited constantly in both layers:
\begin{itemize}[nosep]
  \item \texttt{oid\_in\_stack\_implies\_not\_in\_heap} / \texttt{oid\_in\_heap\_implies\_not\_in\_stack}:
    from \textit{L1}, an id found by one search isn't found by the other.
  \item \texttt{oid\_loc\_rgn\_iff\_in\_heap}, \texttt{oid\_loc\_stk\_iff\_in\_stack},
    \texttt{rid\_loc\_rgn\_iff\_in\_heap}: rewrite $\mathsf{OId}(oid).loc(\langle S,H\rangle) =
    \mathsf{Rgn}(rid)$ into $\exists\, region.\ H(rid) = region \wedge oid \inN region.objMap$.
    Every later file reasons about where an id lives via one of these, not the raw definition.
\end{itemize}

\subsection{$RuntimeConfig.start$ satisfies \textit{ValidConfig}}

One root region (Open, bridge object $0$), one root frame. Each invariant unfolds to a fact about a
one- or two-element list, closed by \texttt{decide}/\texttt{simp}/\texttt{List.nodup\_singleton}.
Exists only to anchor the induction below.

\subsection{The preservation proof pattern}
\label{sec:preservation-pattern}

Each \file{Gc/Model/Mutation/<Op>.lean} defines the operation as an $Option\ RuntimeConfig$
\texttt{do}-block, each bind a failable precondition matching \file{main.tex}'s rule, paired with
an \op{<op>\_cases} lemma unpacking a success into an explicit disjunction, one disjunct per rule
variant. This is the operation's real specification: a missing disjunct or wrong side condition
wouldn't error here (the case lemma just replays the \texttt{do}-block), only later, as an
unprovable or vacuous theorem downstream. Every proof below destructs
$op\ cfg = \mathsf{some}\ cfg'$ via \op{<op>\_cases} once, then argues from the values it produces.

\file{Preservation/<Op>.lean} proves ten lemmas $\op{<op>\_L1}, \dots, \op{<op>\_HS2}$, combined
into $\op{<op>\_valid} : \mathit{ValidConfig}(cfg) \to op\ cfg = \mathsf{some}\ cfg' \to
\mathit{ValidConfig}(cfg')$. \file{Preservation.lean} aggregates all nine, case-split on the
reified $Stmt$ type, into one $\mathit{AllPreserve}(\mathit{ValidConfig})$.

\section{How the proofs actually work}
\label{sec:howtheywork}

Five walkthroughs, covering both layers and every distinct technique. \S\ref{wt:merge-valid} needs
only \S\ref{sec:invariants}. \S\ref{wt:confined}--\ref{wt:merge-portal} need
\S\ref{sec:reach-recap}'s recap of \file{Gc/Reachability/Reachable/} (full definitions in
\file{main.tex}).

\subsection{\op{merge} keeps the configuration valid}
\label{wt:merge-valid}

\textbf{Claim.} \op{merge x} folds a Closed region \op{rid'} into the current Open region \op{rid}
(combines their \textit{objMap}s, deletes \op{rid'}, rebinds \texttt{x} to \op{rid'}'s old bridge
object). Result is still a \textit{ValidConfig}: no duplicate ids (\textit{L1}), contents still
resolve into their own region (\textit{H3}).

\textbf{Intuition.} Pouring one region's objects into another creates and destroys nothing, so ids
stay unique and everything still points home. Right story, but ``nothing created or destroyed'' is
a semantic claim the proof must earn from \op{merge}'s precondition, not its description.

\textbf{The two gaps.}
\begin{itemize}[nosep]
  \item \emph{Disjointness isn't given.} Nothing says the two regions' object-id sets are
    disjoint. If not, the union could identify two different objects (breaking \textit{L1} the
    other way) or, since $AList.union$ erases duplicate keys from one side, silently drop an
    object. The precondition only gives Open/Closed; since those are mutually exclusive and
    \textit{L1} forbids an id in two heap entries, that alone forces disjointness.
  \item \emph{\textit{H3} needs a fresh search.} An object from the deleted \op{rid'} must resolve
    into \op{rid} post-merge, but $.loc(\cdot)$ has no memory: it searches whatever heap exists
    now. The proof reconstructs, via disjointness again, that this search lands on \op{rid} alone.
\end{itemize}

\textbf{The argument.}
\begin{enumerate}[nosep]
  \item \texttt{merge\_corollary\_rid\_ne\_regionId}: \op{rid}, \op{rid'} distinct, since Open
    $\ne$ Closed.
  \item \texttt{merge\_corollary\_region\_unique} (\textit{L1}): no id in two heap entries, so the
    regions share none.
  \item \texttt{merge\_corollary\_union\_append}: on disjoint keys, $s_1 \cup s_2$'s entries are
    just $s_1.entries \mathbin{+\!\!+} s_2.entries$, not $AList.union$'s general, lossier form.
  \item \texttt{merge\_corollary\_heap\_entries\_perm}: post-merge heap entries are a permutation
    of the merged pair plus everything untouched (two nested \texttt{List.exists\_of\_kerase}
    extractions).
  \item \textit{L1} then follows from permutation-invariance of \texttt{List.count}/\texttt{Nodup}.
  \item \textit{H3}: \texttt{merge\_corollary\_loc\_of\_mem}, since the stack side of
    $.loc(\cdot)$'s search is untouched (\op{merge} only rewrites the current frame's varMap) and
    the heap side's uniqueness pins \op{rid}'s merged entry as the only match.
\end{enumerate}

The hard part was never the union: it was showing the two regions could never have shared an id,
then rebuilding where every affected object resolves under the new heap.

\subsection{Reachability, recapped}
\label{sec:reach-recap}

$a \to b$ ($\mathit{ReachableStep}$): one hop. Beyond following an object's stored references,
$\mathsf{RId}(rid) \to b$ needs $H(rid)$ Closed and forces $b = \mathsf{OId}(H(rid).bridgeObjectId)$.
A Closed region's $\mathsf{RId}$ has exactly one outgoing step; an Open region's has none
(\textit{L2}: a step forces Closed, contradicting an on-stack owner). A chain steps into a
suspended region through this one door, never out, which is why this layer exists and why
\S\ref{wt:cr6} is non-trivial.

$\mathit{RegionReachable}(rid, \cdot)$/$\mathit{FrameReachable}(fid, \cdot)$: inductive closures of
$\to$ from a region's bridge object, or a frame's roots ($\mathit{FrameRoot}$: a variable's value,
or its region's bridge object). $\mathit{StackReachable}$: $\mathit{FrameReachable}$ from some
on-stack frame. \file{main.tex} restates all as ``a root, plus a
\texttt{Relation.ReflTransGen}$(\to)$ chain,'' the form every proof below inducts on.
\texttt{ReflTransGen} has two induction principles: constructor induction (fixes the first element,
grows the last; traces forward from a known root) and \texttt{head\_induction\_on} (fixes the last,
peels off the first; chases an escaping chain from a known start). The wrong one doesn't error, it
leaves a stuck goal.

\subsection{Reaching backward in time}
\label{wt:confined}

\textbf{Claim (\texttt{region\_container\_confined}, \texttt{stack\_container\_confined}).} A
reference chain from frame $fid$ back to an object inside an Open region (or stack frame) implies
that container's owner has index $\le fid$.

\textbf{Intuition.} You can't reach an earlier frame's stuff from a later one unless something
already linked the two, close to \textit{S2}/\textit{S3} already.

\textbf{What's missing.} \textit{S2}/\textit{S3} bound one hop. The claim is about a chain of
unknown length, an opaque \texttt{ReflTransGen} witness. Going from one hop to the whole chain
needs induction re-establishing the bound at every hop, with a case split at the moment a chain
crosses from region to stack.

\textbf{The argument.}
\begin{enumerate}[nosep]
  \item Restate as $\mathit{ReflTransGen}$ (\S\ref{sec:reach-recap}); the inductive relation's own
    recursion runs the wrong way.
  \item Induct forward from the fixed start, generalized over the chain's end. Inside the region,
    \texttt{predecessor\_of\_region\_object} (\textit{H3}: one hop from an in-region object stays
    in-region or escapes to stack, never a bare $\mathsf{RId}$ or another region) lets you recurse.
    On the stack, the one-hop bound (\texttt{stack\_step\_region\_index\_le}, \textit{S3} at
    single-hop level) closes it directly.
  \item The two ways a chain starts (a variable slot, a frame's bridge object) close via
    \textit{S3} plus \texttt{merge\_corollary\_regionId\_unique\_index} (same region $\Rightarrow$
    same index).
\end{enumerate}

Here ``can't reach backward in time'' stops being a slogan: the chain is unboundedly long, so a
one-hop fact needs bootstrapping into an arbitrary-length one, with a real case split at the
region/stack boundary.

\subsection{What ``Closed'' really buys you}
\label{wt:cr6}

\textbf{Claim (\textit{CR6}).} Once Closed: if nothing outside holds the region's $\mathsf{RId}$,
nothing inside is stack-reachable; if something does, stack-reachability inside equals
region-reachability from the bridge object.

\textbf{Intuition.} Closed means sealed, front door or nothing, close to the definition already.

\textbf{Why it isn't just the definition.} $\mathit{ReachableStep}$ only gives the forward
direction (one successor). It says nothing about predecessors, so nothing rules out some other way
in already. ``Sealed'' is a negative claim, and proving a negative means closing every other
entrance: another region's object holding a pointer in (ruled out by \textit{H3}: a region's stored
references resolve back into it, so nothing external holds a live pointer in), and a stack frame
still executing there (ruled out by \textit{L2}: an on-stack frame can't own a Closed region).

\textbf{The argument.}
\begin{enumerate}[nosep]
  \item \texttt{predecessor\_of\_closed\_region\_object}: \S\ref{wt:confined}'s lemma, Closed
    version. ``Escaped to stack'' is now impossible (\textit{L2}'s contrapositive,
    \texttt{closed\_region\_not\_owned}); the $\mathsf{RId}$ branch (the portal) becomes available
    instead of ruled out.
  \item \texttt{RegionReachable\_of\_FrameReachable\_closed} / \\
    \texttt{FrameReachable\_rid\_of\_closed\_container}: backward induction, dispatching on (1)'s
    two cases, still inside (recurse) or portal just used (stop; the tail already witnesses the
    $\mathsf{RId}$ was reachable).
  \item \textit{CR6}: stack-reachability inside a Closed region factors uniquely through the
    portal.
\end{enumerate}

The definition of Closed gives the one allowed way in; \textit{CR6} proves it's the only way in,
needing \textit{H3} and \textit{L2} to cover the two entrances the definition alone leaves open.

\subsection{A field write in the active frame}
\label{wt:fieldasgn}

\textbf{Claim ($\mathit{FRALF}$ preserved by \op{fieldAsgn}).} An object in an earlier (possibly
suspended) frame's region, already visible from every later frame, stays visible after a field
write in the active frame.

\textbf{Intuition.} A suspended region's contents aren't touched by a write elsewhere, almost the
whole truth, which is exactly the danger: it stops one level too early.

\textbf{Where it needs correcting.} The write creates one new edge, and nothing about ``active''
rules out that edge connecting two things not connected before. The claim holds because anything
visible from a strictly-suspended frame is, by \S\ref{wt:confined}, already confined to a root no
later than that frame, strictly before the write. So the new edge is provably out of reach of any
chain that matters here.

\textbf{The argument.}
\begin{enumerate}[nosep]
  \item Let $oid_C$ be the mutated object's id. By \S\ref{wt:confined} applied to $oid_C$, anything
    reaching it comes from index $\ge fid_{active}$, the active frame's own.
  \item Any chain confined below $fid_{active}$ never touches $oid_C$ and transports untouched,
    covering every suspended frame automatically.
  \item The remaining case: a chain rooted at the active frame. Walk it forward
    (\texttt{Relation.ReflTransGen.head\_induction\_on}), checking each hop against the new one. A
    match jumps onto the write's new value, which has its own root (\texttt{resolveV\_frameRoot});
    the argument restarts there, back in step 2's case.
\end{enumerate}

The new edge really could matter; the proof shows only one frame's chain could route through it,
and even that reduces to an ordinary confined argument once traced to its own root.

\subsection{Deleting a region reference}
\label{wt:merge-portal}

\textbf{Claim ($\mathit{FRALF}$ preserved by \op{merge}).} Same property, for the operation that
deletes region \op{rid'} and the reference $\mathsf{RId}(rid')$ naming it.

\textbf{Intuition.} \op{merge} relabels a heap key, creates nothing, so reachability from an
earlier frame shouldn't change.

\textbf{What that leaves out.} $\mathsf{RId}(rid')$, in the merging frame's variable \texttt{x}, is
deleted and replaced by \op{region'.bridgeObjectId}. Any argument routed through the old reference
needs re-deriving through the new one; ``nothing used it'' has to be shown, not asserted.

\textbf{The argument.}
\begin{enumerate}[nosep]
  \item \texttt{merge\_ridPrime\_not\_step\_target}: $\mathsf{RId}(rid')$ can't be a step's target,
    or \textit{H2}'s one-reference bound breaks (\texttt{x}'s occurrence is already the allowed
    one). Pure counting, no chain analysis.
  \item So it's only ever a chain's root: ``used the deleted reference'' is a question about the
    start, not a hop-by-hop trace.
  \item \texttt{merge\_chain\_transport\_down}: any chain rooted elsewhere transports
    unconditionally.
  \item The remaining case, rooted at $\mathsf{RId}(rid')$: \texttt{merge\_ridPrime\_reflTransGen\_iff}
    shows its only first hop is its own bridge object, so the chain is the same, one hop shorter,
    rooted at \op{region'.bridgeObjectId}, exactly \texttt{x}'s new value.
\end{enumerate}

\op{merge} needs its own argument because it deletes rather than edits a reference:
\textit{H2}'s bound shows the deleted reference could only be a chain's start, turning an
open-ended re-routing problem into a substitution.

\section{Reference: everything else}

The remaining operations, lemma catalog, and pitfalls, for lookup.

\subsection{The other operations}

\begin{description}
  \item[\op{varAsgn x yf}] Bridge var: reassigns the region's bridge object; \textit{H1} holds
    since the rule's own side condition ($yfRef.loc(\cdot) = \mathsf{Rgn}(rid)$, $rid =
    frame.regionId$) forces the new bridge into the same, unchanged $objMap$. Otherwise: plain
    local write gated by $resolveV(\texttt{x}, \cdot) = \epsilon$, keeping $resolveV$ unambiguous.
  \item[\op{fieldAsgn xf y}] The stack/region rule split shows why \textit{H3} needs
    $status = Open$: a region write needs the writer's frame in that region and the new value
    resolving there too. Drop either and \textit{H3} fails. Reachability proof: \S\ref{wt:fieldasgn}.
  \item[\op{swap x yf}] Busiest rule syntactically, but its invariant proofs are just
    \op{fieldAsgn}'s and \op{varAsgn}'s run side by side.
  \item[\op{makeObjStack x} / \op{makeObjRegion x}] Fresh empty object at
    $\langle S,H\rangle.freshObjectId$. \textit{L1}: \texttt{freshObjectId\_not\_mem}.
    \textit{H2}/\textit{HS1}: vacuous. \op{makeObjRegion} needs $status = Open$.
  \item[\op{makeRegion x}] New Closed region, fresh bridge object. \textit{H2} vacuous by
    freshness; Closed keeps \textit{L2} trivial.
  \item[\op{merge x}] \S\ref{wt:merge-valid}, \S\ref{wt:merge-portal}.
  \item[\op{enter xf bridgeVar} / \op{exit}] Push/pop, flip status. \textit{L2}/\textit{S1}: index
    bookkeeping. Pushed frame is empty, so \textit{S2}/\textit{S3}/\textit{HS1} add nothing.
    Reachability: ``additive'' (\S\ref{sec:fralf-shapes}), \texttt{safe\_reflTransGen\_transport}
    applies directly.
\end{description}

\subsection{$\mathit{FRALF}$ preservation: three shapes}
\label{sec:fralf-shapes}

\begin{enumerate}[nosep]
  \item \textbf{Additive} (\op{enter}, \op{exit}, \op{makeObjStack}, \op{makeObjRegion},
    \op{makeRegion}, \op{varAsgn}): content unchanged, \texttt{safe\_reflTransGen\_transport}
    (\S\ref{sec:catalog}) moves any chain for free.
  \item \textbf{Content-mutating} (\op{fieldAsgn}, \op{swap}): \S\ref{wt:fieldasgn}'s escape
    argument, owner-index bound.
  \item \textbf{Reference-deleting} (\op{merge}): \S\ref{wt:merge-portal}'s root-only argument,
    \textit{H2} count bound.
\end{enumerate}
\file{StackReachableInvariant/} proves the same case analysis as an $\iff$. The two directions
aren't mirror images: the new edge exists in $cfg'$ not $cfg$, so escape only arises going
$cfg' \to cfg$.

\subsection{Reusable lemma catalog}
\label{sec:catalog}

\begin{description}
  \item[\texttt{stackWithIndex\_frame\_transport\_down/up\_of\_shape\_eq}] $cfg,cfg'$ agreeing up
    to projection $\pi$ at every stack position transports $stackWithIndex$ membership, preserving
    $.index$ and $\pi$.
  \item[\texttt{stackWithIndex\_objMap\_get\_eq\_of\_last\_varMap\_update}] Last frame's $varMap$
    changed $\Rightarrow$ $objMap$ unaffected everywhere.
  \item[\texttt{merge\_corollary\_regionId\_unique\_index}] \textit{S1} at the $FrameWithIndex$
    level: same region $\Rightarrow$ same frame $\Rightarrow$ same index.
  \item[\texttt{swap\_corollary\_stackWithIndex\_index\_inj} / \texttt{\_find\_eq}] Same
    $.index$ $\Rightarrow$ same element. The most-cited lemma in \file{Reachable/}: collapses
    differently-named-but-identical frame variables.
  \item[\texttt{predecessor\_of\_region\_object} / \texttt{\_of\_closed\_region\_object}] Backward
    confinement, one hop; \S\ref{wt:confined}/\S\ref{wt:cr6}.
  \item[\texttt{SafeRef} / \texttt{predecessor\_safe} / \texttt{safe\_reflTransGen\_transport}] Safe
    for $rid$: in $rid$ or on stack, never a bare $\mathsf{RId}$. Propagates backward one hop;
    matching step relations transport a safe-ending chain wholesale. Behind the ``additive'' shape.
  \item[\texttt{stack\_container\_confined} / \texttt{region\_container\_confined}]
    \S\ref{wt:confined}'s index bound.
  \item[\texttt{RegionReachable\_of\_FrameReachable\_closed} / \\
    \texttt{FrameReachable\_rid\_of\_closed\_container}] Closed-region confinement behind
    \textit{CR6}.
  \item[\textit{H2}-forbids-a-step-target trick] (\S\ref{wt:merge-portal}) generalizes: deleting
    any uniquely-referenced region reference reduces transport to a root-only case split.
\end{description}

\subsection{Pitfalls}

\begin{pitfall}[\texttt{AList.insert} reorders on an existing key]
Via $kinsert$/$kerase$, unlike a fresh key. Entries-level equality after an insert
(\S\ref{wt:merge-valid}) needs \texttt{List.Perm} (\texttt{exists\_of\_kerase} +
\texttt{NodupKeys.eq\_of\_mk\_mem}), not ``insert = append.''
\end{pitfall}

\begin{pitfall}[\textit{Status} has no \texttt{LawfulBEq}]
\texttt{beq\_iff\_eq} doesn't fire. Open/Closed contradictions close by case-splitting
$region.status$ and \texttt{decide}, not rewriting \texttt{==}.
\end{pitfall}

\begin{pitfall}[the \op{<op>\_cases} lemma is the semantics]
Read it, not the opaque \texttt{do}-block, when in doubt about a precondition. A mistake there
stays invisible until a downstream theorem quietly stops being provable.
\end{pitfall}

\begin{pitfall}[Open-only tools on a Closed region]
\S\ref{wt:confined}'s region lemma needs $status = Open$ (leans on
\texttt{predecessor\_of\_region\_object}, Open-only). For Closed, use \S\ref{wt:cr6}'s lemmas: the
Open-only one just fails to apply, an easy mismatch when adapting to a new operation.
\end{pitfall}

\begin{pitfall}[$\mathit{loc}(\cdot)$ treated as opaque]
Unfolded once, in \file{Validity.lean}, to prove \S\ref{sec:invariants}'s three
\texttt{\_iff\_in\_*} lemmas; every later file treats it as opaque from then on.
\end{pitfall}

\begin{pitfall}[wrong \texttt{ReflTransGen} induction principle]
Constructor induction fixes the start, grows the end; \texttt{head\_induction\_on} fixes the end,
peels the start. Wrong choice: a stuck goal, not an error. See \S\ref{sec:reach-recap}.
\end{pitfall}

\begin{pitfall}[an $\iff$ needs both halves independently]
$cfg' \to cfg$ isn't $cfg \to cfg'$ flipped: the new edge exists only in $cfg'$, so escape arises
in one direction only. See \S\ref{sec:fralf-shapes}.
\end{pitfall}

\subsection{What's deliberately not done}
No \op{Guarantees.lean}-style top-level GC-safety theorem yet. Everything above is single-step;
composing over an arbitrary-length trace (the literal report.pdf claim) isn't currently pursued.

\end{document}
